\documentclass[10pt,a4paper]{amsart}
\usepackage[margin=19mm]{geometry}
\usepackage{amsmath,amssymb,mathtools}
\usepackage[scr=rsfs,cal=euler]{mathalfa}
\usepackage{xfrac}
\usepackage[unicode,hidelinks,bookmarks=false]{hyperref}
\newcommand{\ie}{i.e.\ }
\newcommand{\cf}{cf.\ }
\newcommand{\resp}{respectively\ }
\newcommand{\Subsection}{\subsection}
\providecommand{\nobreakdash}{\nobreak\hbox{-}}
\setlength{\emergencystretch}{2em}
\multlinegap0pt
\usepackage{amssymb}
\usepackage{tikz}
\usepackage{cancel}
\usepackage{enumerate}
\usepackage{mathtools}
\usepackage{float}
\usepackage{dsfont}
\usepackage[shortlabels]{enumitem}
\newtheorem{theorem}{Theorem}
\newtheorem{remark}{Remark}
\newtheorem{lemma}{Lemma}
\def\R{\mathbb{R}}
\def\cC{\mathcal{C}}
\def\cN{\mathcal{N}}
\def\cR{\mathcal{R}}
\def\cS{\mathcal{S}}
\def\k{\kappa}
\def\la{\lambda}
\title{\vspace*{-1.4cm}Identifiability of SDEs for reaction networks}
\author{Louis Faul, Linard Hoessly, Panqiu Xia}
\date{Source: arXiv:2505.07638v3 (CC BY 4.0)}
\begin{document}
\maketitle

\noindent\emph{Source and licence.} \vspace*{-1.4cm}Identifiability of SDEs for reaction networks. Version arXiv:2505.07638v3 (CC BY 4.0), distributed by arXiv under the Creative Commons Attribution 4.0 International licence, http://creativecommons.org/licenses/by/4.0/, as recorded in the arXiv licence metadata for this exact version and shown on the abstract page https://arxiv.org/abs/2505.07638v3. Full licence notice: \texttt{licenses/arxiv-2505.07638-CC-BY-4.0.txt}.\par\smallskip

\noindent\textit{Editorial scope.} Counted unit: the article's theorem theorem1 (source0.tex lines 765--777). The article's complete original proof of it is delivered with it (source0.tex lines 939--972, one \texttt{proof} environment of 34 lines, which the article places away from the statement). No mathematics is written, repaired or re-derived: the statement and the proof are the article's own lines, in the article's own notation, and no retained line is substituted or re-flowed.  Retained uncounted as the source-local context the proof needs: lines 351--353; lines 355--358; lines 398--400; lines 718--720; lines 780--783; lines 925--925; lines 928--934.  Nothing was omitted from any retained span, so the package carries no omission record.  No retained line contains a citation, so the wrapper carries no bibliography and no bibliography entry.  No retained line contains a figure environment, an image-inclusion command or drawing code, so no image is delivered and the package declares no asset dependency.  Editorial formatting only, in addition: an AMS \texttt{amsart} preamble that restates the article's own declarations and macros used by the retained lines, namely the environments \texttt{theorem}, \texttt{remark}, \texttt{lemma} on separate counters, together with the standard packages those lines need.  The retained environments print in this document as Theorem 1, Remark 1, Lemma 1 in order of appearance, whereas the article prints the same items with the numbering of its own full document.  The complete native source of the article as distributed in the arXiv e-print for this exact version is bundled unchanged as \texttt{sources/177971/source0.tex}.
\begin{equation}\label{drift_vector}
      A_i(x) = \sum_{r=1}^d \Gamma_{i,r}\la_r(x) = \sum_{r=1}^d (y'_{i,r} - y_{i,r}) \kappa_r x^{y_r}, \quad i = 1,\dots, n;
\end{equation}

\begin{equation}\label{diffusion_matrix}
      B_{i,j}(x)=
(\sigma(x) \sigma(x)^{\intercal})_{i,j} = \sum_{r=1}^d \Gamma_{ir}\Gamma_{jr} \la_r(x) = \sum_{r = 1}^d (y'_{i,r} - y_{i,r}) (y'_{j,r} - y_{j,r}) \kappa_r x^{y_r}, \quad i, j = 1, \dots, n.
\end{equation}

\begin{equation}\label{generalized_reaction_vectors}
    \left\{\left(y'-y,(y' - y)\cdot (y' - y)^{\intercal}\right) \big| \, y\to y'\in\cR\right\}
\end{equation}

\begin{equation}\label{generator}
    Lf(x) = \langle A(x), \nabla \rangle f{(x)}+{\frac{1}{2}} (\nabla^{\intercal} \cdot B (x) \cdot \nabla) f{(x)} = \sum_{i = 1}^n A_i(x)\frac{\partial}{\partial x_i}f(x)+\frac{1}{2}\sum_{i,j = 1}^n B_{i,j}(x)\frac{\partial^2}{\partial x_i\partial x_j}f(x),
\end{equation}

\begin{remark}\label{rmk-theorem1}
   If two RNs $\cN$ and $\cN'$ are unconfoundable w.r.t. their generators, then an analogous of property
(2) in Theorem \ref{theorem1} holds for all rate constant vectors $\kappa \in \mathbb{R}^d_{>0}$ and $\kappa' \in \mathbb{R}^{d'}_{>0}$; see Subsection \ref{prf_theorem1}.
\end{remark}

\subsection{Proof of Theorem \ref{theorem1}} \label{prf_theorem1} 

\begin{lemma}\label{lem:monomial_independence}
Let $Y\subset \mathbb Z_{\ge 0}^n$ be a finite set and let $v_y\in\mathbb R^m$ be vectors for $y\in Y$. If
\[
\sum_{y\in Y} v_y\,x^y = 0 \qquad \text{for all } x\in U,\quad U\subseteq\mathbb R^n\text{ open},
\]
then $v_y=0$ for all $y\in Y$.
\end{lemma}

\begin{theorem}\label{theorem1}
Let $\cN = (\cS, \cC, \cR)$ be an RN. The followings are equivalent.
\begin{enumerate}
\item $\cN$ is reaction-identifiable w.r.t. its generator.

\item There is a non-empty bounded open set $U\subseteq\R^n_{>0}$, such that for all rate constant vector{s} $\kappa \neq \kappa'\in\R^d_{>0}$ the corresponding generators are different on $C^{\infty}_c(U)$.

\item For every source complex $y\in \cC$, the set of vectors in \eqref{generalized_reaction_vectors}
is linearly independent.

\end{enumerate}

\end{theorem}

\begin{proof}[Proof of Theorem~\ref{theorem1}]$(1)\implies (2)$: Let $\cN$ be an RN such that $(2)$ fails. Then, there exist rate constant vectors $\kappa \neq \kappa'$, whose corresponding generators coincide on a nonempty open set $U$. Recall that the generators take the form \eqref{generator}. The fact that two generators coincide on $C_c^{\infty} (U)$ is equivalent to their coefficients $A_i$ and $B_i$, which are polynomials, coinciding on $U$. Since a nonzero polynomial cannot vanish on a nonempty open set, it follows that the coefficients (and hence the generators) coincide on all of $\R^n_{>0}$ (resp. on $C^\infty_c(\mathbb{R}^n_{>0})$), see Lemma \ref{lem:monomial_independence}. As a consequence, $\cN$ is not identifiable w.r.t. its generator; that is $(1)$ fails. This proves the implication $(1) \implies (2)$.

Notice that in the above argument, we do not actually require the underlying RNs with rate constant vectors $\kappa$ and $\kappa'$ to share the same network structure. Therefore, the proof still justifies Remark \ref{rmk-theorem1} without any modification.


$(2) \implies (3)$: Assume that on an open subset $U$ as given the generators are different for all $\kappa \neq \kappa'$, but that there is a source complex $y \in \cC$ such that the set of vectors 
$$\left\{\left(y' - y,(y' - y)\cdot (y' - y)^{\intercal}\right) \big| \, y\to y'\in\cR\right\}$$
is linearly dependent. Then there exist some real numbers $\alpha_{y\to y'}$, not all zero  such that
\begin{align*}
\left(\sum_{y' \colon y \to y'\in\cR} \alpha_{y\to y'}(y' - y),\sum_{y' \colon y \to y'\in\cR} \alpha_{y\to y'} (y' - y)\cdot (y' - y)^{\intercal} \right)  
= 0.
\end{align*}
Hence we can choose rate constant vectors $\k,\k' \in \R^{d}_{>0}$ such that $\k_{y \to y'}- \k'_{y \to y'} = \alpha_{y \to y'}$ for all coefficients of reactions $y \to y'\in\cR$, and $\k_{R} = \k'_{R} = 1$ for other reactions $R \in\cR$. Then, recalling \eqref{drift_vector}, \eqref{diffusion_matrix}, and \eqref{generator}, for these rate constant vectors,
with $L$ and $L'$ denoting the generator with $k$ and $k'$ respectively,
\begin{align*}
 L f(x) - L' f(x) = & \sum_{y\to y'\in \cR} \left\langle (y' - y)(\k_{y\to y'}-\k'_{y\to y'})x^y, \nabla\right\rangle f \\ 
& + \frac{1}{2} \sum_{y\to y'\in \cR} \left(\nabla^{\intercal} \cdot (y' - y)\cdot(y' - y)^{\intercal}(\k_{y\to y'}-\k'_{y\to y'}) x^y \cdot \nabla \right) f  = 0.
\end{align*}
This contradicts with our assumption, and thus verifies $(2) \implies (3)$.

$(3) \implies (1)$: 
 Assume that for all source complexes $y \in \cC$ the set of vectors 
$$\left\{\left(y' - y,(y' - y)\cdot (y' - y)^{\intercal}\right) \big| \, y\to y'\in\cR \right\}$$
is linearly independent. Let $\k,\k'\in\R^d_{>0}$ such that $L f(x) - L' f(x) = 0$, for all test functions $f\in C^{\infty}_c(\R_{> 0}^n)$ and $x \in R_{> 0}^n$. In other words,

$$\sum_{y \in \cC} \left(\sum_{y' \colon y\to y'\in \cR}(y' - y)(\k_{y\to y'} - \k'_{y\to y'})x^y,\sum_{y' \colon y\to y'\in \cR}(y' - y)\cdot(y' - y)^{\intercal}(\k_{y\to y'}-k'_{y\to y'})x^y \right)= 0,$$
where the right hand side is the zero-function $\R^n\to \R^{n}\times\R^n \otimes \R^n$. This yields that
$$\sum_{y\in \cC}x^y\sum_{y' \colon  y\to y'\in \cR} (\k_{y\to y'}-\k'_{y\to y'}) \left((y' - y),(y' - y)\cdot(y' - y)^{\intercal}\right)= 0.$$
Since these vectors are polynomials in $x$ vanishing on $\R^n_{>0}$, their polynomial coefficients must be zero. Therefore, for each source complex $y\in\cC$,
$$
\sum_{y' \colon y\to y'\in \cR}(\k_{y\to y'}-\k'_{y\to y'}) \left((y' - y),(y' - y)\cdot(y' - y)^{\intercal}\right)= 0 .
$$
By the linear independence, we get that $\k_{y\to y'}-\k'_{y\to y'}=0$. This concludes the proof of $(3) \implies (1)$. The proof of this theorem is complete.
\end{proof}

\end{document}
